\subsubsection{Disjoint Sparse Table}
\begin{lstlisting}[style=mystyle, language=C++]
// TC: O(N log N) build, O(1) query
// usa: queries idempotentes en arreglos estaticos (alternativa a Sparse Table)
#include <bits/stdc++.h>
using namespace std;

struct DisjointSparseTable{
	int n, LOG;
	vector<vector<int>> st;
	vector<int> lg;
	DisjointSparseTable(const vector<int>& a){
		n = a.size();
		LOG = 32 - __builtin_clz(n + 1);
		st.assign(LOG, vector<int>(n));
		st[0] = a;
		int len = 1;
		for(int k=1;k<LOG;k++){
			for(int i=0 ; i<n ; i++){
				int j = i + len;
				if(j >= n){ st[k][i] = st[k-1][i]; continue; }
				st[k][i] = min(st[k-1][i], st[k-1][j]);
			}
			for(int i=n-1 ; i>=0 ; i--){
				int j = i - len;
				if(j < 0){ 
					st[k][i] = (i+1 < n ? min(st[k][i], st[k-1][i+1]) : st[k-1][i]); 
					continue; 
				}
				st[k][i] = min(st[k-1][i], st[k-1][j]);
			}
			len <<= 1;
		}
		lg.assign(n + 1, 0);
		for(int i=2 ; i<=n ; i++) lg[i] = lg[i/2] + 1;
	}
	int query(int l, int r){
		if(l == r) return st[0][l];
		int k = lg[l ^ r];
		return min(st[k][l], st[k][r]);
	}
};

int main(){
	DisjointSparseTable dst({3, 1, 4, 1, 5, 9, 2, 6});
	cout << dst.query(0, 4) << "\n";  // 1
	cout << dst.query(2, 7) << "\n";  // 2
	return 0;
}
\end{lstlisting}
